\documentclass[10pt,a4paper]{amsart}
\usepackage{amsmath,amssymb,amsthm,graphicx}
\usepackage{verbatim}
\newenvironment{reference}{\comment}{\endcomment}
%\newenvironment{reference}{}{}
\newenvironment{slogan}{\comment}{\endcomment}
\newenvironment{history}{\comment}{\endcomment}

% For commutative diagrams we use Xy-pic
\usepackage[all]{xy}

% We use 2cell for 2-commutative diagrams.
\xyoption{2cell}
% Theorem environments.
%
\theoremstyle{plain}
\newtheorem{theorem}[subsection]{Theorem}
\newtheorem{proposition}[subsection]{Proposition}
\newtheorem{lemma}[subsection]{Lemma}

\theoremstyle{definition}
\newtheorem{definition}[subsection]{Definition}
\newtheorem{example}[subsection]{Example}
\newtheorem{exercise}[subsection]{Exercise}
\newtheorem{situation}[subsection]{Situation}

\theoremstyle{remark}
\newtheorem{remark}[subsection]{Remark}
\newtheorem{remarks}[subsection]{Remarks}

\numberwithin{equation}{subsection}

% Macros
%
\def\lim{\mathop{\mathrm{lim}}\nolimits}
\def\colim{\mathop{\mathrm{colim}}\nolimits}
\def\Spec{\mathop{\mathrm{Spec}}}
\def\Hom{\mathop{\mathrm{Hom}}\nolimits}
\def\Ext{\mathop{\mathrm{Ext}}\nolimits}
\def\SheafHom{\mathop{\mathcal{H}\!\mathit{om}}\nolimits}
\def\SheafExt{\mathop{\mathcal{E}\!\mathit{xt}}\nolimits}
\def\Sch{\mathit{Sch}}
\def\Mor{\mathop{\mathrm{Mor}}\nolimits}
\def\Ob{\mathop{\mathrm{Ob}}\nolimits}
\def\Sh{\mathop{\mathit{Sh}}\nolimits}
\def\NL{\mathop{N\!L}\nolimits}
\def\CH{\mathop{\mathrm{CH}}\nolimits}
\def\proetale{{pro\text{-}\acute{e}tale}}
\def\etale{{\acute{e}tale}}
\def\QCoh{\mathit{QCoh}}
\def\Ker{\mathop{\mathrm{Ker}}}
\def\Im{\mathop{\mathrm{Im}}}
\def\Coker{\mathop{\mathrm{Coker}}}
\def\Coim{\mathop{\mathrm{Coim}}}

% Boxtimes
%
\DeclareMathSymbol{\boxtimes}{\mathbin}{AMSa}{"02}

%
% Macros for moduli stacks/spaces
%
\def\QCohstack{\mathcal{QC}\!\mathit{oh}}
\def\Cohstack{\mathcal{C}\!\mathit{oh}}
\def\Spacesstack{\mathcal{S}\!\mathit{paces}}
\def\Quotfunctor{\mathrm{Quot}}
\def\Hilbfunctor{\mathrm{Hilb}}
\def\Curvesstack{\mathcal{C}\!\mathit{urves}}
\def\Polarizedstack{\mathcal{P}\!\mathit{olarized}}
\def\Complexesstack{\mathcal{C}\!\mathit{omplexes}}
% \Pic is the operator that assigns to X its picard group, usage \Pic(X)
% \Picardstack_{X/B} denotes the Picard stack of X over B
% \Picardfunctor_{X/B} denotes the Picard functor of X over B
\def\Pic{\mathop{\mathrm{Pic}}\nolimits}
\def\Picardstack{\mathcal{P}\!\mathit{ic}}
\def\Picardfunctor{\mathrm{Pic}}
\def\Deformationcategory{\mathcal{D}\!\mathit{ef}}
\usepackage{cleveref}
\crefname{originalequation}{original equation}{original equations}
\Crefname{originalequation}{Original equation}{Original equations}
\makeatletter
\newlabel{equation-localize-ringed}{{04LZ}{source}{}{sourcepointer-1}{}}
\newlabel{equation-localize-ringed@cref}{{[originalequation][0][]04LZ}{[1][0][]source}{}{}{}}
\makeatother
\crefname{originalreference}{original reference}{original references}
\Crefname{originalreference}{Original reference}{Original references}
\makeatletter
\newlabel{sites-modules-lemma-relocalize-morphism-ringed-sites}{{04J1}{source}{}{sourcepointer-2}{}}
\newlabel{sites-modules-lemma-relocalize-morphism-ringed-sites@cref}{{[originalreference][0][]04J1}{[1][0][]source}{}{}{}}
\makeatother
\crefname{originallemma}{original lemma}{original lemmas}
\Crefname{originallemma}{Original lemma}{Original lemmas}
\makeatletter
\newlabel{lemma-etale-permanence}{{03FV}{source}{}{sourcepointer-3}{}}
\newlabel{lemma-etale-permanence@cref}{{[originallemma][0][]03FV}{[1][0][]source}{}{}{}}
\makeatother
\crefname{originalequation}{original equation}{original equations}
\Crefname{originalequation}{Original equation}{Original equations}
\makeatletter
\newlabel{equation-localize}{{04LY}{source}{}{sourcepointer-4}{}}
\newlabel{equation-localize@cref}{{[originalequation][0][]04LY}{[1][0][]source}{}{}{}}
\makeatother
\crefname{originalreference}{original reference}{original references}
\Crefname{originalreference}{Original reference}{Original references}
\makeatletter
\newlabel{sites-lemma-localize-morphism}{{03CF}{source}{}{sourcepointer-5}{}}
\newlabel{sites-lemma-localize-morphism@cref}{{[originalreference][0][]03CF}{[1][0][]source}{}{}{}}
\makeatother
\title{Localization compatibility of morphisms of ringed topoi over etale covers}
\author{The Stacks Project Authors}
\date{stacks-project:spaces-properties.tex:local-snapshot-20261012; GNU FDL 1.2 or later}
\begin{document}
\maketitle
\noindent\textit{Editorial note.} All delivered mathematical text is the author's own, copied verbatim from the pinned Stacks Project chapter file spaces-properties.tex; no statement or proof step was rewritten, reordered or removed.\par\medskip
\noindent\textit{Editorial note.} Span recomputation by label pairing (LANE\_BRIEF 10.2): in this chapter the author's proof environment follows the statement environment as a standalone block; the spec covers the complete environment pair (statement span, proof span) as recorded in the pool candidate.\par\medskip
\noindent\textit{Editorial note.} Occupation scan: combined claim\_key + statement\_anchor + generic-line scan over all registration files (active and withdrawn) found 13 S10-geom published cards for spaces-properties.tex (prop-occupied.json, 2026-10-14); those 13 labels are excluded from this batch by design. For the 77 free labels: 0 occupied. The registry mentions field (7557 files) and focused-seeds (645) both point to the same S10-geom cards (dual-source check), no other occupation.\par\medskip
\ifdefined\hypertarget\hypertarget{sourcepointer-1}{}\fi
\noindent\textit{Source reference.} Original equation 04LZ: sources/stacks-chow/tags.txt, source lines 6025--6025. Cross-chapter standard prerequisite cited by the author inside the counted span; it is used as a pointer and its own proof is not part of this item.\par\medskip
\ifdefined\hypertarget\hypertarget{sourcepointer-2}{}\fi
\noindent\textit{Source reference.} Original reference 04J1: sources/stacks-chow/tags.txt, source lines 5922--5922. Cross-chapter standard prerequisite cited by the author inside the counted span; it is used as a pointer and its own proof is not part of this item.\par\medskip
\ifdefined\hypertarget\hypertarget{sourcepointer-3}{}\fi
\noindent\textit{Source reference.} Original lemma 03FV: sources/stacks-chow/tags.txt, source lines 4506--4506. Cross-chapter standard prerequisite cited by the author inside the counted span; it is used as a pointer and its own proof is not part of this item.\par\medskip
\ifdefined\hypertarget\hypertarget{sourcepointer-4}{}\fi
\noindent\textit{Source reference.} Original equation 04LY: sources/stacks-chow/tags.txt, source lines 6024--6024. Cross-chapter standard prerequisite cited by the author inside the counted span; it is used as a pointer and its own proof is not part of this item.\par\medskip
\ifdefined\hypertarget\hypertarget{sourcepointer-5}{}\fi
\noindent\textit{Source reference.} Original reference 03CF: sources/stacks-chow/tags.txt, source lines 4382--4382. Cross-chapter standard prerequisite cited by the author inside the counted span; it is used as a pointer and its own proof is not part of this item.\par\medskip
\begin{lemma}
\label{lemma-relocalize-morphism}
Let $S$ be a scheme. Let
$$
\xymatrix{
U \ar[d]_p \ar[r]_g & V \ar[d]^q \\
X \ar[r]^f & Y
}
$$
be a commutative diagram of algebraic spaces over $S$ with $p$ and $q$ \'etale.
Via the identifications
(\ref{equation-localize-ringed}) for $U \to X$ and $V \to Y$
the morphism of ringed topoi
$$
(g_{spaces, \etale}, g^\sharp) :
(\Sh(U_{spaces, \etale}), \mathcal{O}_U)
\longrightarrow
(\Sh(V_{spaces, \etale}), \mathcal{O}_V)
$$
is $2$-isomorphic to the morphism $(f_{spaces, \etale, c}, f_c^\sharp)$
constructed in
Modules on Sites,
Lemma \ref{sites-modules-lemma-relocalize-morphism-ringed-sites}
starting with the morphism of ringed sites
$(f_{spaces, \etale}, f^\sharp)$ and
the map $c : U \to V \times_Y X$ corresponding to $g$.
\end{lemma}

\begin{proof}
The morphism $(f_{spaces, \etale, c}, f_c^\sharp)$ is defined as a
composition $f' \circ j$
of a localization and a base change map. Similarly $g$ is a composition
$U \to V \times_Y X \to V$. Hence it suffices to prove
the lemma in the following two cases: (1) $f = \text{id}$, and
(2) $U = X \times_Y V$. In case (1) the morphism $g : U \to V$ is
\'etale, see
Lemma \ref{lemma-etale-permanence}.
Hence $(g_{spaces, \etale}, g^\sharp)$ is a localization morphism
by the discussion surrounding
Equations (\ref{equation-localize}) and
(\ref{equation-localize-ringed})
which is exactly the content of the lemma in this case.
In case (2) the morphism $g_{spaces, \etale}$
comes from the morphism of ringed sites given by the functor
$V_{spaces, \etale} \to U_{spaces, \etale}$,
$V'/V \mapsto V' \times_V U/U$
which is also what the morphism $f'$ is defined by, see
Sites, Lemma \ref{sites-lemma-localize-morphism}.
We omit the verification that $(f')^\sharp = g^\sharp$
in this case (both are the restriction of $f^\sharp$
to $U_{spaces, \etale}$).
\end{proof}

\end{document}
